% ======================================================================
% Diagrama de flujo — generado automáticamente por generar_diagrama_flujo.py
% NO EDITAR MANUALMENTE; regenerar desde el descriptor JSON.
% ======================================================================
\documentclass[11pt,a4paper]{article}

% ── Codificación, fuentes y lenguaje ───────────────────────────────────────────
\usepackage[utf8]{inputenc}
\usepackage[T1]{fontenc}
\usepackage[default,scale=0.95]{sourcesanspro}
\renewcommand{\familydefault}{\sfdefault}
\usepackage[spanish,es-noshorthands,es-tabla]{babel}

% ── Geometría y espaciado ─────────────────────────────────────────────────────
\usepackage[top=2.2cm, bottom=2.5cm, left=2.2cm, right=2.2cm, headheight=15pt]{geometry}
\usepackage{setspace}
\setstretch{1.18}
\usepackage{parskip}

% ── TikZ (simbolos ISO 5807 / ANSI X3.5) ──────────────────────────────────────
\usepackage{tikz}
\usetikzlibrary{babel, arrows.meta, positioning, shapes.geometric,
                shapes.symbols, decorations.pathmorphing, calc}

% ── Tablas y listas ───────────────────────────────────────────────────────────
\usepackage{booktabs}
\usepackage{tabularx}
\usepackage{array}
\usepackage{enumitem}

% ── Colores, cajas e iconos ───────────────────────────────────────────────────
\usepackage[dvipsnames]{xcolor}
\usepackage{tcolorbox}
\tcbuselibrary{skins, breakable}
\usepackage{fontawesome5}

% ── Secciones con barra de color (infografía) ─────────────────────────────────
\usepackage{titlesec}
\titleformat{\section}
  {\normalfont\LARGE\bfseries\color{azulNoche}}
  {\colorbox{cyanNeon}{\color{fondoOscuro}\thesection}}{0.7em}{}
  [\vspace{2pt}\color{cyanNeon}\rule{\linewidth}{1.8pt}]
\titlespacing*{\section}{0pt}{24pt}{10pt}
\titleformat{\subsection}
  {\normalfont\large\bfseries\color{azulNoche}}
  {\textcolor{cyanNeon}{\thesubsection}}{0.7em}{}
  [\vspace{1pt}\color{grisLinea}\rule{\linewidth}{0.6pt}]
\titlespacing*{\subsection}{0pt}{14pt}{6pt}

% ── Cabeceras y pies de página ────────────────────────────────────────────────
\usepackage{fancyhdr}
\usepackage{lastpage}
\pagestyle{fancy}
\fancyhf{}
\renewcommand{\headrulewidth}{0pt}
\renewcommand{\footrulewidth}{0pt}
\fancyfoot[C]{%
  \begin{minipage}{\linewidth}
    \vspace{2pt}
    \color{cyanNeon}\rule{\linewidth}{1.2pt}\\[2pt]
    \centering\color{grisTexto}\footnotesize
    Página \thepage\ de \pageref{LastPage}
  \end{minipage}%
}

% ── Hiperenlaces ──────────────────────────────────────────────────────────────
\usepackage[colorlinks=true, linkcolor=azulNoche, urlcolor=cyanNeon]{hyperref}
\sloppy
\emergencystretch 3em

% ── Paleta de colores ─────────────────────────────────────────────────────────
\definecolor{fondoOscuro}{HTML}{0D1B2A}
\definecolor{verdeTurquesa}{HTML}{004D40}
\definecolor{azulNoche}{HTML}{1B3A6B}
\definecolor{azulMedio}{HTML}{2A5298}
\definecolor{cyanNeon}{HTML}{00C8E0}
\definecolor{verdeSignal}{HTML}{00B887}
\definecolor{naranjaVivo}{HTML}{FF7043}
\definecolor{rojoAlerta}{HTML}{E53935}
\definecolor{grisPapel}{HTML}{F4F6FA}
\definecolor{grisLinea}{HTML}{C8D0E0}
\definecolor{grisTexto}{HTML}{6B7A99}
\definecolor{amarilloNota}{HTML}{FFB300}

% ── Tarjetas ──────────────────────────────────────────────────────────────────
\tcbset{
  cajaContenido/.style={
    enhanced, breakable,
    arc=5pt, outer arc=5pt,
    colback=grisPapel, colframe=grisLinea,
    boxrule=0.8pt,
    fonttitle=\bfseries\small, coltitle=azulNoche,
    top=6pt, bottom=6pt, left=10pt, right=10pt,
  },
}

% ── Macros ────────────────────────────────────────────────────────────────────
\newcommand{\iconotexto}[2]{\textcolor{cyanNeon}{\faIcon{#1}}~#2}
\newcommand{\iconoBanda}{sitemap}
\newcommand{\bandaTitulo}[3][fondoOscuro]{%
  \noindent
  \begin{tcolorbox}[%
    enhanced, arc=4mm, colback=#1!10!white, colframe=#1, boxrule=0pt,
    borderline west={6pt}{0pt}{cyanNeon},
    top=6pt, bottom=6pt, left=8pt, right=8pt,
  ]
    \noindent
    \begin{minipage}[c]{0.10\linewidth}
      \centering\color{cyanNeon}\faIcon{tachometer-alt}
    \end{minipage}%
    \hspace{8pt}%
    \begin{minipage}[c]{0.88\linewidth}
      {\fontsize{20}{24}\selectfont\bfseries\color{white}#2}\par\vspace{5pt}%
      \textcolor{cyanNeon!80!white}{\small\faIcon{angle-right}~#3}%
    \end{minipage}
    \vspace{4pt}
  \end{tcolorbox}%
  \vspace{8pt}%
}


  % ── Estilos de flujo (ISO 5807 / ANSI X3.5) ─────────────────────────────
  \tikzset{
    flujo/.style={
      >=Stealth,
      font=\footnotesize\sffamily,
    },
    terminador/.style={
      draw=azulNoche, fill=azulNoche!8!white, rounded corners=3pt,
      text width=1.6cm, align=center, inner sep=2mm,
      font=\footnotesize\sffamily,
    },
    proceso/.style={
      draw=azulNoche, fill=cyanNeon!10!white,
      text width=1.6cm, align=center, inner sep=2mm,
      font=\footnotesize\sffamily,
    },
    entradasalida/.style={
      draw=azulNoche, fill=verdeSignal!10!white,
      trapezium, trapezium left angle=75, trapezium right angle=105,
      text width=1.8cm, align=center, inner sep=2mm,
      font=\footnotesize\sffamily,
    },
    decision/.style={
      draw=azulNoche, fill=amarilloNota!20!white,
      diamond, aspect=1.6, text width=1.5cm, align=center, inner sep=1pt,
      font=\footnotesize\sffamily,
    },
    almacenamiento/.style={
      draw=azulNoche, fill=grisPapel,
      cylinder, shape border rotate=90, aspect=0.3,
      text width=1.7cm, align=center, inner sep=2mm,
      font=\footnotesize\sffamily,
    },
    documento/.style={
      draw=azulNoche, fill=azulMedio!8!white,
      text width=1.8cm, align=center, inner sep=2mm,
      font=\footnotesize\sffamily,
    },
    nota/.style={
      draw=grisLinea, dashed, fill=grisPapel,
      text width=1.8cm, align=center, inner sep=2mm,
      font=\scriptsize\sffamily,
    },
    etiqueta/.style={
      font=\scriptsize\sffamily, fill=white,
      fill opacity=0.75, text opacity=1, inner sep=1pt,
    },
  }


% ── Cabecera del documento ────────────────────────────────────────────────────
\fancyhead[L]{\small\color{azulNoche}\textbf{ELECTRONICA --- capa 2 (orquestacion)}}
\fancyhead[R]{\small\color{grisTexto}flujo\_auditar\_sistema.py}
\renewcommand{\headrulewidth}{0pt}

% ─────────────────────────────────────────────────────────────────────────────
\begin{document}

\bandaTitulo{Auditoria de sistema}{RAM, CPU, servicios y GPU: medir, veredictar y no tocar nada}

\vspace{4pt}
\begin{center}
\begin{tcolorbox}[cajaContenido, title={\faIcon{map-signs}~Leyenda de símbolos---ISO 5807 / ANSI X3.5}]
\small
Óvalo = \textbf{terminador} (inicio/fin) $\cdot$
Rectángulo = \textbf{proceso} (acción determinista) $\cdot$
Rombo = \textbf{decisión} (ramas Sí/No) $\cdot$
Paralelogramo = \textbf{entrada/salida} (lectura/escritura de archivos) $\cdot$
Cilindro = \textbf{almacenamiento online} (estado, snapshots) $\cdot$
Rectángulo con base ondulada = \textbf{documento} (sesiones, logs) $\cdot$
Rectángulo punteado gris = \textbf{nota explicativa} (sin acción).
Flujo: arriba$\rightarrow$abajo, izquierda$\rightarrow$derecha.

\faIcon{tags}~Cada símbolo lleva una \textbf{etiqueta} (T/P/D/E/A/N + número, según su tipo);
su significado se expande en la tabla \emph{Leyenda de etiquetas} bajo cada diagrama.
\end{tcolorbox}
\end{center}

\vspace{6pt}

\section{\iconotexto{cogs}{La pasada: de la peticion al codigo de salida}}

\begin{center}
\begin{tikzpicture}[flujo, >=Stealth, x=1cm, y=1cm]
  % Fondo decorativo
  \fill[azulNoche!3!white, rounded corners=4pt]
    (-2.3,-0.4) rectangle (5.8,-22.1);
  \node[terminador] (T1) at (0.0,-2.2) {T1};
  \node[proceso] (P1) at (0.0,-4.4) {P1};
  \node[proceso] (P2) at (0.0,-6.6) {P2};
  \node[proceso] (P3) at (0.0,-8.8) {P3};
  \node[decision] (D1) at (0.0,-11.0) {D1};
  \node[proceso] (P4) at (0.0,-13.2) {P4};
  \node[decision] (D2) at (0.0,-15.4) {D2};
  \node[almacenamiento] (A1) at (0.0,-17.6) {A1};
  \node[terminador] (T2) at (0.0,-19.8) {T2};
  \node[nota] (N1) at (3.5,-6.6) {N1};
  \node[nota] (N2) at (3.5,-11.0) {N2};
  \node[nota] (N3) at (3.5,-15.4) {N3};

  \draw[->] (T1.south) -- (P1.north);
  \draw[->] (P1.south) -- (P2.north);
  \draw[->] (P2.south) -- (P3.north);
  \draw[->] (P3.south) -- (D1.north);
  \draw[->] (D1.south) -- (P4.north);
  \draw[->] (D1.east) -- (N2.west);
  \draw[->] (P4.south) -- (D2.north);
  \draw[->] (D2.south) -- (A1.north);
  \draw[->] (A1.south) -- (T2.north);
  \draw[->] (P2.east) -- (N1.west);
  \draw[->] (D2.east) -- (N3.west);
\end{tikzpicture}
\end{center}

\vspace{6pt}
\begin{tcolorbox}[cajaContenido, title={\faIcon{table}~Leyenda de etiquetas}]
\small
\begin{tabularx}{\linewidth}{@{}>{\centering\arraybackslash}p{1.4cm}X@{}}
\toprule
\textbf{Etiqueta} & \textbf{Significado} \\ \midrule
A1 & Persistir el informe y el estado en .tmp/, y trazar inicio, cada dimension y fin en el log append-only. Si la traza falla, el flujo avisa y continua. \\
D1 & La medicion es completa o ha quedado un hueco? Un hueco no se rellena con una suposicion: se declara no\_verificado. \\
D2 & Cual es el PEOR estado presente? Gana el maximo, nunca un promedio ni una media ponderada. \\
N1 & El algebra de veredictos se comparte con flujo\_auditar\_repo.py desde execution/veredicto\_algebra.py. Duplicarlo produciria dos verdades divergiendo en silencio. \\
N2 & La capa 3 devuelve hechos y la capa 2 decide. Un test falla si una medicion anade la clave estado. \\
N3 & El informe incluye los umbrales aplicados: la misma maquina con otros umbrales daria otro veredicto. \\
P1 & Acotar: resolver el subconjunto de las 5 dimensiones y validarlo. Un nombre desconocido sale con 3, reservado al uso incorrecto de flags. \\
P2 & Medir con la capa 3: leer /proc/meminfo, /proc/swaps, /proc/loadavg, /proc/stat, /proc/pressure y /sys/block/zram*, y consultar systemctl show. Solo hechos, jamas un estado. \\
P3 & Interpretar con funciones puras (medicion) -> dimension, aplicando los umbrales. Sin ficheros, sin reloj, sin procesos: por eso se testean enteras. \\
P4 & Componer el informe con ensamblar\_informe, otra funcion pura, e imprimir la tabla legible con lo que hay que hacer en orden de gravedad. \\
T1 & Se pide auditar el sistema, con las dimensiones opcionales acotadas por --dimension o --solo. \\
T2 & Salir con el codigo del veredicto: 0 limpio o con avisos, 1 con fallos, 2 sin verificar, 3 uso incorrecto de flags. \\
\bottomrule
\end{tabularx}
\end{tcolorbox}


\section{\iconotexto{sticky-note}{Notas}}
\begin{itemize}
\item El veredicto es el PEOR estado presente, nunca un promedio.
\item Un hueco de medicion impide el verde: sale con 2, nunca con 0.
\item Los umbrales viajan dentro del informe; sin ellos el veredicto es irreproducible.
\item Este flujo no purga, no reinicia y no necesita sudo. Purgar es otro flujo.
\end{itemize}

% ─────────────────────────────────────────────────────────────────────────────
\end{document}
